


\subsection{Ring-LPN overview and its use in PCGs} \label{subsec:ringlpn}
The Ring-LPN assumption works over a quotient ring $\mathbb R:=\mathbb F_q[X]/F(X)$ (degree $N$), with “sparse” noise polynomials of Hamming weight $t$. Writing $e,f\gets \mathsf{HW}_{\mathbb R,t}$ for $t$-sparse polynomials and $r\gets \mathbb R$, the distribution
\[
(r,\; r\cdot e+f)\qquad\text{is computationally indistinguishable from}\qquad (r,\; u),\ u\gets \mathbb R,
\]
as formalized in \cite[Def.~3.1]{boyle2020efficient}. Concretely, $\mathbb R$ will be the polynomials modulo $F(X)=(X^N+1)$. This choice allows efficient polynomial multiplications via a negacyclic NTT.

Boyle et al.\ showed how to realize a \emph{silent} PCG for OLE (and authenticated triples) from Ring-LPN by packing many field OLEs into one ring instance and expanding with fast polynomial arithmetic; see their overview and the OLE/triples sections \cite{boyle2020efficient}. In brief:
\begin{itemize}
  \item Party $p\in \{0,1\}$ samples two (or more) random $t$-sparse polynomials $(e_{p,0}, e_{p,1})\in \mathbb R^2$. Concretely, this will represent sampling $t$-sparse vectors. 
  \item The parties will sample a random polynomial $r\in \mathbb R$.
  \item Party $p$ will compute $x_p' := e_{p,0} + r \cdot e_{p,1}$.
  \item Party $p$ applies an NTT to obtain $x_p:=\mathsf{NTT}(x_p')\in \mathbb{F}_q^N$. A vector of $N$ $\mathbb F _q$ elements.
  \item The parties \emph{want} to compute a secret sharing 
  \begin{align*}
    \share{y'}&=x_0' \cdot x_1'\in \mathbb R\\
             &= ( e_{0,0} + r \cdot e_{0,1}) \cdot ( e_{1,0} + r \cdot e_{1,1})\\
             &= 
             e_{0,0} \cdot e_{1,0}  + 
             r \cdot e_{0,1} \cdot  e_{1,0} + 
             e_{0,0}\cdot r \cdot e_{1,1} +
             r^2 \cdot  e_{0,1} \cdot e_{1,1}
  \end{align*}
  \item Given $\share{y'}$, the parties can then locally compute $\share{y}:=\mathsf{NTT}(\share{y'})$. The critical property is that
  $$
        x_{0} \odot x_{1} = \share{y}_0 + \share{y_1}
  $$
  where $\odot$ is componentwise multiplication over vectors $\mathbb{F}_q^N$. This is percisely an Oblivious Linear Evaluation (OLE) correlation. One interpretation of this transform is that the NTT is performing polynomails evaluation. If  $\share{y'}=x_0' \cdot x_1'\in \mathbb R$, then for some $\alpha\in \mathbb F_q$ the same holds for $\share{y'(\alpha)}=x_0'(\alpha) \cdot x_1'(\alpha)\in \mathbb F_q$ (for an appropriately chosen $\mathbb R$).    Two such correlations can then be locally converted into $N$ Beaver triples
  $$
    \share{a},\share{b}, \share{c} \in \mathbb{F}_q^N \text{  s.t. } \share{a} \odot \share{b} = \share{c}
  $$
  
  \item Let $A_{p,i}\subset [0,N)^t, B_{p,i}\in \mathbb F_q^t$ denote the positions and values of the non-zeros coefficients in $e_{p,i}\in \mathbb R$. Observe that the product polynomial $e_{0,i}\cdot e_{1,j}$ has position and values described as
  \begin{align*}
        A'_{i,j}:=&A_{0,i} \boxplus A_{1,j} \in [0,2N)^{t^2},\\
        B'_{i,j}:=&B_{0,i} \otimes B_{1,j} \in \mathbb{F}_q^{t^2}
  \end{align*}
  where $\boxplus$ is the cartisian addition and $\otimes$ is the cartisian product of the sets. That is, 
    $$
        e_{0,i}\cdot e_{1,j} = \mathsf{makePoly}(A'_{i,j}, B'_{i,j}) \mod F(X)
    $$
    where $\mathsf{makePoly}(P\in \mathbb N^t, V\in \mathbb F^t)= \sum_i V_i X^{P_i}$ returns the polynomails with the $V$ coefficients at the $P$ positons.

  \item Observe that ${\mathsf{makePoly}(\share{A'_{i,j}}, \share{B'_{i,j}})}$ can directly be constructed from a DMPF. In particular, ${\mathsf{makePoly}(\share{A'_{i,j}}, \share{B'_{i,j}})}$ is perciesely the same function as a DMPF when the polynomial is express as a vector of coefficients. 
\end{itemize}

\noindent
The outline above gives a concrete protocol for generating $N$ OLE and Beaver triples. 
\begin{enumerate} 
\item The parties sample their sparse polynomails. 
\item Each party construct the sparse represetation of their polynomails. Using generic MPC or similar, they compute the cartisian addition and product of the positons and values to obtains $\share{A'_{i,j}}, \share{B'_{i,j}}$. 
\item They use a DPMF of size $t^2$ for each $(i,j)$  to obtain $\mathsf{makePoly}(A'_{i,j}, B'_{i,j})$ in secret shared form. 
\item They reduce this modulo $F(X)$.
\item They perform local additions and FFT operations to obtain $N$ OLE or $N/2$ Beaver Triples over $\mathbb F_q$. 
\end{enumerate}


One can further optimize this protocol by restricting the polynomials to being regular\cite{boyle2020efficient}. Given this restriction, $\mathsf{makePoly}$ can be implemented as $t$ instances of a DMPF, each of size $t$. This can allow for a more efficient implement than a single DMPF instance of size $t^2.$





\subsection{Stationary LPN / SSD and amortizing noise generation}\label{sec:stationary}

\emph{Stationary LPN} (a.k.a.\ Stationary Syndrome Decoding, SSD) captures the setting where the \emph{support} of the noise is reused across many instances while the \emph{payloads} (coefficients) are freshly sampled each time \cite{Stationary}.
In our Ring-LPN notation from above, fix for each party $p\in\{0,1\}$ and selector $i\in\{0,1\}$ a support set
\[
L_{p,i}\subseteq[0,N),\qquad |L_{p,i}|=t,
\]
and let $A_{p,i}$ be the sorted list of positions in $L_{p,i}$.
For the $s$-th expansion, sample fresh coefficient vectors $B^{(s)}_{p,i}\in\mathbb F_q^{t}$ and form
\[
e^{(s)}_{p,i}\ :=\ \mathsf{makePoly}\!\big(A_{p,i},\,B^{(s)}_{p,i}\big)\ \in\ \mathbb R,
\]
while sampling an independent $r^{(s)}\leftarrow\mathbb R$.
Security asserts that reusing $(L_{p,i})$ across many $s=1,2,\ldots,q$ (with fresh coefficients and ring masks) remains pseudorandom against the relevant classes of attacks, so that the pairs
\(
\big(r^{(s)},\,r^{(s)}\cdot e^{(s)}_{p,1}+e^{(s)}_{p,0}\big)
\)
are computationally indistinguishable from random in~$\mathbb R^2$ for each~$s$.

\paragraph{Consequence for the product terms.}
With stationary supports, the \emph{product supports} are fixed once and for all:
\[
A'_{i,j}\ :=\ A_{0,i}\ \boxplus\ A_{1,j}\ \in [0,2N)^{t^2},\qquad
B^{(s)}_{i,j}\ :=\ B^{(s)}_{0,i}\ \otimes\ B^{(s)}_{1,j}\ \in \mathbb F_q^{t^2},
\]
and
\(
e^{(s)}_{0,i}\cdot e^{(s)}_{1,j}=\mathsf{makePoly}\!\big(A'_{i,j},\,B^{(s)}_{i,j}\big)\bmod F(X).
\)
Thus, across $s=1,\dots,q$, the \emph{addresses} $A'_{i,j}$ never change; only the \emph{values} $B^{(s)}_{i,j}$ do.

\paragraph{How Stationary applies to our pipeline.}
The Ring-LPN PCG flow in \sectionref{subsec:ringlpn} already factors the heavy work into:
(i) building (and then multiplying) the sparse polynomials, and
(ii) routing those sparse contributions into the ring (and then into the NTT domain).
Stationarity lets us push nearly all position-dependent work \emph{out of the per-expansion loop}:

\begin{enumerate}
  \item \textbf{One-time (setup).}
  \begin{itemize}
    \item Fix supports $A_{p,0},A_{p,1}$ and derive the four fixed product supports $A'_{i,j}$.
    \item Precompute the position-only routing for each $A'_{i,j}$ with our new DMPF constructions: this fixes the level tags, permutations, and correction structure once. No payloads are embedded yet.
  \end{itemize}

  \item \textbf{Per expansion $s=1,\dots,q$.}
  \begin{itemize}
    \item Sample fresh coefficients $B^{(s)}_{p,i}$ (and fresh ring mask $r^{(s)}$); form $B^{(s)}_{i,j}=B^{(s)}_{0,i}\otimes B^{(s)}_{1,j}$.
    \item \emph{Value-only programming:} reuse the cached routing for each fixed $A'_{i,j}$ and inject the new $B^{(s)}_{i,j}$ values (no recomputation of tags or inner-product gadgets). This is exactly the “payload injection” stage in our DMPF box, repeated with new shares.
    \item Reduce $\sum_{i,j}\mathsf{makePoly}(A'_{i,j},B^{(s)}_{i,j})$ modulo $F(X)$ and proceed with the NTT and componentwise products to obtain the packed OLEs / triples.
  \end{itemize}
\end{enumerate}

\paragraph{Cost implications.}
Let $D=\lceil\log_2 N\rceil$ and recall $|A'_{i,j}|=t^2$.
\begin{itemize}
  \item \emph{Fresh (non-stationary):} each expansion repeats position discovery and routing, rebuilding the  structure for every $(i,j)$, plus programming values—this is where the original distributed approach paid an $n$-dependent price via $\SharedIPD$.
  \item \emph{Stationary:} we pay the position-routing cost \emph{once}. Each expansion programs only the new values $B^{(s)}_{i,j}$ through the cached DPMF skeleton and then runs ring/NTT arithmetic. In our improved DPMF path, the per-expansion nonlinear work is $O(\mathrm{poly}(t))$ for the DMPF value path, independent of $N$; the arithmetic dominates (NTT and pointwise multiplies).
\end{itemize}

\paragraph{Practical takeaway.}
For degree-2 correlations (Ring-LPN $\to$ packed OLE / triples), Stationary LPN perfectly matches the factorization in our construction:
\begin{center}
\emph{fix supports once} \(\Rightarrow\) \emph{precompute DPMF routing once} \(\Rightarrow\) \emph{per expansion: only new coefficients and ring arithmetic.}
\end{center}
This reduces the online cryptographic work to value injection on a fixed sparse footprint and lets the NTT-backed Goldilocks arithmetic drive overall throughput, which is precisely the regime where our implementation attains multi-million ops/sec.
